\chapter{Conclusion}\label{ch:conclusion}
This thesis addressed the challenge of evaluating and comparing \ac{rl} agents for power grid operation. Motivated by the increasing complexity of modern power grids and the growing interest in using \ac{rl} agents to support grid operators, the research question was formulated: \begin{displayquote}
How can a benchmarking tool be designed to quantitatively and qualitatively evaluate and compare agents for power grid operation, providing both an overall performance score and detailed feedback on their strengths and weaknesses across different operating scenarios?
\end{displayquote}
To answer this question, the benchmark \benchmark was designed, implemented and evaluated through a series of experiments, presented in Chapters \ref{ch:methodology}, \ref{ch:implementation} and \ref{ch:experiments}. This chapter concludes the thesis. Section \ref{sec:contributions-findings} summarises the main contributions of \benchmark and the key findings of the conducted experiments, relating them back to the research question. Section \ref{sec:future-work} then outlines possible directions for future work that emerged during the development and evaluation of the benchmark.




\section{Contributions and Findings}\label{sec:contributions-findings}
This section summarises the main contributions of this thesis and relates them back to the research question. Section \ref{ssec:contribution-benchmark} presents the contributions made through the design and implementation of \benchmark, outlining how its features address the research question from a design perspective. Section \ref{ssec:contribution-experiments} summarises the key findings of the experiments conducted in Chapter \ref{ch:experiments}, validating the benchmark's capabilities empirically. Finally, Section \ref{ssec:contribution-takeaway} condenses the most important insight of this thesis into a single, concise takeaway.

\subsection{Benchmark}\label{ssec:contribution-benchmark}
The developed benchmark \benchmark, presented in this work, is a configurable metrics- and scenario-based benchmark for agents working in power grid operation on the transmission grid. Each scenario specifies one power grid, provided by the IEEE test cases or the SimBench dataset, and one profile from the SimBench dataset, where specific indices from the profile are chosen for evaluation. The grid of a scenario is modifiable, with different possible grid modifications that are configurable by the user. Grid modifications create extreme scenarios with imbalanced grid components. This includes increased generation or consumption in one generator or load, maintenance events that make grid components unusable for the scenario, or capacity scaling, which limits the possible usage of lines or transformers. 

For benchmarking agents, \benchmark offers different options. Either one of two baseline agents can be chosen, or a custom agent created by the user. To use a custom agent, two blueprint classes exist, from which the custom agent inherits. The \texttt{BenchAgent} defines the necessary interfaces that \benchmark needs to interact with a custom agent, such as agent setup, training and action selection. The \texttt{RayAgent} is a subclass of the \texttt{BenchAgent} and offers an interface for agents that use the Ray framework for building their \ac{rl} architecture. It simplifies agent definition for Ray agents, as only the agent initialisation needs to be implemented. Furthermore, the \texttt{BaseTuner} exists for hyperparameter tuning of an agent before training. This class also has a subclass, the \texttt{OptunaTuner}, which simplifies making hyperparameter tuners for custom agents using the Optuna optimisation framework. For each scenario, a custom agent can be tuned using a customised tuning class, and is then trained on the scenario using user-specified training indices from the scenario profile. After training, the evaluation indices from the scenario profile are each executed individually to evaluate the agent.

The evaluation of each run in a scenario is built on a multi-dimensional metric framework. It measures multiple grid features, and converts them into a normalised score, each within the same range of 0 to 100. To summarise all the gathered scores, the scores are aggregated and form a hierarchical scoring system. Each layer of the scoring system summarises the performance of the grouped scores underlying the layer. The category score summarises multiple metric scores, using configurable weights for each metric. The run score summarises multiple category scores, with each category containing a configurable weight as well. Finally, the scenario score summarises the run scores, and the total score summarises all scenario scores, on average. For fixed run data and deterministic agent behaviour, the benchmark evaluation and scoring procedure is reproducible. Full experimental reproducibility additionally depends on controlling the agent's stochasticity and random seeds.

The gathered results can be presented using different included tools. These create statistics, visualisations with plots and web apps, and comparisons between agents. With the help of these, the benchmark results are summarised and presented in a way that the user can analyse every layer of the hierarchical scoring system and how well the agent performed for specific scenarios or metrics. In addition to these tools, the replay tool creates a visualisation from the run data gathered during the scenario execution. The replay tool offers an analysis of the actions taken by an agent and how stable the grid was in each timestep. This provides another analytical layer to the benchmark, where users can understand the agent's behaviour and scoring.

The multi-dimensional metrics and hierarchical scoring system make this benchmark different from the past works presented in Section \ref{sec:rw-benchmarking}. Furthermore, it has more diagnostic tools built in, which make a stronger analysis into the agent's performance possible. With regard to the research question, the benchmark presented in this work addresses both its quantitative and qualitative aspects. It provides a quantitative evaluation using normalised metric scoring, offering a variety of metrics to measure and evaluate different grid features. These metric scores are then summarised using a hierarchical aggregation system, which gives an overall performance score. Furthermore, it delivers a qualitative evaluation by offering visual and textual reporting formats and a replay tool, which reveal how an agent behaves, not just how well it scores. The output saved by the benchmark does not just store the aggregated total score. It also provides information on each layer of the benchmark, from specific metric scores in single runs, up to aggregated scenario scores. Additionally, the reporting tools demonstrate the strengths and weaknesses that each agent has for specific metrics, categories or operating scenarios. The scenario tuple design enables diverse, reproducible, and configurable test conditions. 


\subsection{Experiments}\label{ssec:contribution-experiments}
Another important part of this work is the experiments presented in Chapter \ref{ch:experiments}. The experiments tested \benchmark on a variety of scenarios. The experiments are grouped into two categories. One category contains three experiments, evaluating performance measured on the grid operation. The other category contains one experiment, which evaluated computational efficiency.

The three experiments in the grid operation category focussed on the metric categories of overload mitigation, voltage violation mitigation, survival, and costs. Firstly, the standard analysis evaluated how an agent runs with an unmodified grid, and how the benchmark scores that. Next, a diversity analysis tested the functionality to modify power grids and how this affects the agents and benchmark scoring. For this experiment, every type of grid modification that is implemented into the benchmark design was tested, creating a diverse report on how different modifications affect grid complexity. Finally, the scalable modifications were evaluated using a stress test that examined three levels of scaling, each of which decreased grid stability.

For the computational efficiency category, only one experiment was conducted, which only focussed on the metric category of computational performance. This experiment evaluated the same scenarios from the standard analysis, just with a different focus. For each run, the runtimes, CPU and memory usages were measured and used for the computational performance score. This tested how fast and resource-intensive the greedy and \ac{ppo} agent were. The do-nothing agent was excluded from this comparison, since its runtime serves as the baseline against which computational performance is normalised.

The key findings of the experiments are that the \ac{ppo} achieved the highest scores overall. However, it showed how an agent can learn reward hacking, in order to maximise reward, in a way that is undesirable for power grid operation. It also pointed out that no agent is universally best, as the performance of each agent in the experiment was scenario-dependent. Furthermore, the utilised baseline agents provided important context on scenario difficulty. This gives more information on an agent's performance quality. Additionally, the usage of two different grids and many grid modifications, showcased how these aspects affect the difficulty to maintain grid stability. 

The experiments highlighted how important it is for a benchmark to show more than an aggregated score. A high aggregated score does not necessarily imply a desirable operational behaviour for an agent. It is important to retain detailed diagnostic information alongside the overall ranking. Furthermore, the experiments confirm that \benchmark successfully fulfils the research question. While the quantitative scores allowed a clear ranking of agents, the qualitative analysis uncovered critical behavioural flaws, such as the \ac{ppo}'s reward hacking. This information would have remained hidden if only the aggregated total score had been used as a performance measure. This demonstrates the importance of multi-dimensional evaluation for a trustworthy assessment of power grid agents.

\subsection{Takeaway}\label{ssec:contribution-takeaway}
This work demonstrates that the evaluation of power grid agents cannot be reduced to a single number. \benchmark shows that quantitative scoring and qualitative diagnostics must go hand in hand. Otherwise, an agent could achieve a top score while displaying behaviour that would be unacceptable in a real-world grid operation.







\section{Future Work}\label{sec:future-work}
While \benchmark successfully addresses the research question of this thesis, its development and evaluation also revealed several limitations and possible extensions that were beyond the scope of this work. This section presents these directions for future work, grouped into four areas. Section \ref{ssec:future-metrics} discusses possible improvements to the metric and scoring design. Section \ref{ssec:future-scenarios} addresses possible extensions of scenario diversity and realism. Section \ref{ssec:future-safety} presents ideas for incorporating safety constraints and improved failure handling. Finally, Section \ref{ssec:future-agents} discusses future directions regarding agent evaluation and leaderboard functionality.



\subsection{Improved Metric and Scoring Design}\label{ssec:future-metrics}
The experiments presented in Chapter \ref{ch:experiments} revealed several possible improvements to the current metric and scoring design of \benchmark. The following paragraphs discuss four such improvements: incorporating baseline-relative scoring, balancing categories with an uneven number of metrics, extending the costs category, and improving the clipping approach used during metric normalisation.

\paragraph{Baseline-Relative Scoring.}
For the scoring, the baselines can be integrated more. As already mentioned in Section \ref{sec:experiment-learnings}, the baseline agents bring more information, such as scenario difficulty. As some scenarios are much easier or harder than other, treating and evaluating them equally makes comparison with other scenarios less precise. Instead, a future scoring approach could normalise the agent's performance relative to one or more baseline agents, so that the score reflects performance relative to scenario difficulty. This would be similar to the scoring system in \ac{l2rpn}, where a score above 0 means that an agent performs better than a do-nothing agent.

\paragraph{Balanced Categories.}
Furthermore, the categories are unevenly represented in terms of the number of metrics, which affects how granular each category score is. As Table \ref{tab:metric_categories} shows, survival is measured by a single metric, while overload mitigation comprises up to seven metrics that expand into as many as twelve individually weighted components once mean and worst aggregation and $N-1$ variants are included. The categories of costs and voltage violation mitigation each aggregate around seven weighted components. Since a category score is the weighted average of its metrics, as shown in Equation \ref{eq:run_score}, categories with more components benefit from the aggregation. A single poor value has less influence on the score. In survival, however, there is no such smoothing, since its score depends entirely on one measurement. A more balanced category design, whether through additional non-redundant metrics or a scoring mechanism that accounts for metric count directly, is therefore an important addition to \benchmark for future work.

\paragraph{Costs Category.}
For the categories of costs an improvement in metrics would also be beneficial. For example, some mitigation-related metrics may also be cost-related, as exceeding the physical limits of components damages them and increases the costs. For this aspect, metrics measuring equipment stress would be useful.

\paragraph{Clipping Improvement.} 
The current normalisation clips values that are outside a predefined range. 
This can hide exceptionally good performance, or how severely an agent violates a threshold. 
Future scoring could preserve information beyond the clipping range, by either increasing the scoring range or saving this information additionally, and penalising or rewarding the extent to which the clipping range is exceeded. An example of where this would be useful is the metric of maximum line loading. Currently, the score is zero for any line loading over 100\%. However, there is a difference in severity between a line loading of 110\% and one of 200\%. A penalty that depends on the extent to which the safe line-loading percentage is exceeded would be more adequate.



\subsection{Extension of Scenario Diversity and Realism}\label{ssec:future-scenarios}
Besides improvements to the scoring design, the scenario design of \benchmark could also be extended in future work to increase its diversity and realism. The following paragraphs present five possible directions: procedural profile and grid creation, evaluation on renewable energy profiles, broader modifications for larger grids, longer episode lengths, and improved physical realism of the currently supported grids.

\paragraph{Profile and Grid Creation.}
In future work, a profile and grid creator could be useful to create more grid data. This can be used similarly to Procgen, where procedural generation tests agents' generalisation. During the development of \benchmark, there was an experimental feature that made it possible for a user to create a new and challenging power grid through the configurations. However, these power grids still need a lot of improvement and are not yet suitable, as they do not represent realistic power grids. With future work, this kind of feature could be implemented.

\paragraph{Renewable Energy Profiles.}
In the experiments, none of the scenarios had renewable energy sources. As already mentioned in Chapter \ref{ch:background}, these sources are increasingly used and pose a challenge to maintaining grid stability. \benchmark supports profiles with renewable energy sources using the SimBench profiles. The benchmark needs to be tested on these profiles, and the agent's behaviour should be analysed on scenarios with renewable energy sources.

\paragraph{Broader Modifications for Larger Grids.}
In larger grids such as the 30-bus grid, modifying only a small number of components has not been very influential on the scenario difficulty, as the experiments show. In future scenarios, the modifications should scale with grid size and modify several components simultaneously.

\paragraph{Longer Episodes.}
Longer simulations, by increasing the scenarios' episode length, would allow the benchmark to evaluate how long an agent maintains grid stability. The evaluated scenarios in the experiment included only one-day scenarios. However, longer time periods such as those in RL2Grid would be more challenging and would potentially reveal more about the benchmarked agent's performance and behaviour. Also, currently, if a line exceeds its limit, it is put into maintenance for the rest of the episode due to overheating. When the episode length increases, there should be a cooldown time, where the transmission line is repaired and is usable again, within one episode. Also, it would be interesting to test how maintenance events in the middle of the episode would affect grid stability. In the current version of \benchmark, the maintenance events are set for a whole episode.

\paragraph{Physically Realistic Grids.}
As mentioned in Section \ref{sec:discussion-limitations}, the grid in the 14-bus case has some unrealistic physical settings in the transformers and transmission lines. The grid and profile data need to be evaluated in more depth to improve the benchmark quality when these scenarios are chosen.


\subsection{Safety Constraints and Failure Handling}\label{ssec:future-safety}
The reward hacking behaviour observed for the \ac{ppo} agent in Chapter \ref{ch:experiments} highlights the need for stronger safety mechanisms within \benchmark. The following paragraphs discuss four possible directions to address this: introducing penalty metrics, incorporating safety constraints, integrating \ac{tso} operational knowledge, and improving failure diagnostics for non-convergence cases.

\paragraph{Penalty Metrics.}
The actions taken by the \ac{ppo} agent for some scenarios of the experiment should not have received such a high score when most of the grid did not receive any power and most of the power came from the external grid. Metrics should explicitly penalise excessive use of the external grid, islanding, and disconnected loads.

\paragraph{Safety Constraints.}
An even stricter approach against the undesirable strategy of the \ac{ppo} agent would be safety constraints, similar to those in RL2Grid. Instead of a \ac{mdp}, in future work, a \ac{cmdp} would be a better choice for imposing constraints on the possible actions that an agent can choose from. This would remove the possibility for an agent to choose actions, which cause load islanding, excessive load shedding, or disconnecting most of the power grid just to achieve high rewards.

\paragraph{TSO Operational Knowledge.}
Another feature from RL2Grid that is very useful is to incorporate \ac{tso} operational knowledge into the environment. In future work, operator-inspired rules or heuristics could avoid unnecessary actions when the system is already secure enough. These heuristics can then restore the original topology, just like in the RL2Grid benchmark. This shifts the agents' focus more towards unstable grid states that they have to stabilise. 

\paragraph{Non-Convergence Fault Detection.}
In future versions of \benchmark, a diagnostic tool could be created to analyse non-convergence cases. So far, only the pandapower diagnostic tool has been utilised in the experiment. Nevertheless, a structured fault detection in cases of non-convergence would help provide a better understanding of run failures. Currently, the replay functionality also makes such analysis possible. However, sometimes it is not easy to interpret which action or component exactly caused the failure.


\subsection{Agents and Leaderboards}\label{ssec:future-agents}
Finally, future work could also focus on extending the evaluation of agents within \benchmark. The following paragraphs discuss four directions: evaluating a broader range of \ac{rl} architectures, extending the custom reward functionality, studying the effects of benchmark overfitting, and introducing leaderboard functionality to \benchmark.

\paragraph{More Agents.}
As discussed in Section \ref{sec:discussion-limitations}, the number of evaluated agents needs to be increased to expand the coverage of tested \ac{rl} architectures. There are a lot of different agents besides the \ac{ppo} agent, as Section \ref{ssec:example-agents} illustrated. It would be worth evaluating other \ac{rl} architectures such as value-based agents like Q-Learning or \ac{dqn}, or other policy-based agents like \ac{a3c}.

\paragraph{Custom Rewards.}
The current custom reward functionality is limited, as it was not the focus of this work. However, it would be useful to extend the custom reward function by adding more metrics to it. In future work, this could be used to analyse how the reward design influences benchmark performance. For example, the same agent could be tested on different reward functions to study their effect.

\paragraph{Benchmark Overfitting.}
With a fully extended custom reward, another interesting research direction would be the effects of benchmark overfitting. If the user chooses the reward metrics to be as close as possible to the benchmark metrics, it is unclear how this would affect the agents. Would the agents learn to exploit weak spots of the benchmark configuration and scoring, or would they learn a robust grid-operation policy? In the experiments, the \ac{ppo} agent already hinted at a possibility through its reward hacking behaviour. However, a thorough study on this would showcase more examples and, potentially, how this can be avoided.

\paragraph{Leaderboards.}
As \ac{l2rpn} showed, a leaderboard-based experiment provides a significant amount of information on \ac{rl} agents. The same could be done with \benchmark in future work. With a fixed configuration, leaderboards across multiple agents could be implemented. The configurations should contain realistic, validated scenarios, and the leaderboard could be split into groups. A global leaderboard compares the agents' overall performance, while specific leaderboards could evaluate best safety strategies, best cost reduction, or best computational efficiency. Moreover, the leaderboards should not only show the ranking of the best agents, but also display information on other metrics, categories, and scenarios for each agent. This way, multiple agents are comparable across multiple evaluation criteria. This would provide much broader information about the benchmark and how different types of agents perform in it. Furthermore, it would provide a much deeper evaluation of \benchmark's own quality as a benchmarking tool. 